\begin{lemma}\label{lem:bounded-measure-carrier}
Let $\theta\geq0$ be smooth and big, and let $\mu$ be a finite non-pluripolar measure on $X$. There are a bounded $\psi\in\PSH(X,\theta)$ and a nonnegative
$f\in L^1(\mathrm{MA}_\theta(\psi))$ such that
\[
 \mu=f\,\mathrm{MA}_\theta(\psi).
\]

\end{lemma}
\begin{proof}
Apply \cref{prop:relative-measure-carrier} with $\varphi=0$. Semipositivity gives $P_{\epsilon\theta}(0)=0$ for every $0<\epsilon<1$, and $\psi\sim0$ means that $\psi$ is bounded.
\end{proof}

These representations give an $L^1$ density, whereas \cref{thm:bounded-subsolution} requires a bounded density against the measure of a bounded potential. For an unbounded prescribed type, the representing potential is itself unbounded. The general-measure existence argument in \cref{sec:general-measure-relative} uses a bounded Hermitian reference and relative compactness to pass to the equation.
